\documentclass[10pt,twocolumn,letterpaper]{article}

\usepackage{graphicx}
\usepackage{amsmath}
\usepackage{amssymb}
\usepackage{booktabs}
\usepackage{hyperref}
\usepackage{microtype}

\title{QRIntel: A Proactive, Explainable Cybersecurity Intelligence Platform
for Detecting Malicious QR-Based Infrastructure}

\author{
Ashmi KOTIAN\\
Department of Cyber Security, NMAM Institute of Technology (NMAMIT)\\
Nitte (Deemed to be University) $|$ Nitte, Udupi, 574110, Karnataka, India\\
ashmi.kotian@nitte.edu.in
\and
SPOORTHI YADAV\thanks{Corresponding Author}\\
Department of Cyber Security, NMAM Institute of Technology (NMAMIT)\\
Nitte (Deemed to be University) $|$ Nitte, Udupi, 574110, Karnataka, India\\
spoorthipyadav@gmail.com
\and
SINDHOORA R\\
Department of Cyber Security, NMAM Institute of Technology (NMAMIT)\\
Nitte (Deemed to be University) $|$ Nitte, Udupi, 574110, Karnataka, India\\
sindhoora807@gmail.com
}

\date{Received: 1 January 202X $|$ Revised: 2 February 202X $|$ Accepted:
3 March 202X \\
Licensed under a CC-BY 4.0 license $|$ Copyright (c) by the authors
$|$ DOI: https://doi.org/10.48084/etasr.XXXX}

\begin{document}

\maketitle

\begin{abstract}
The extensive use of Quick Response (QR) Codes in mobile banking, retail,
and authentication has made transactions incredibly convenient, yet inherently
vulnerable due to the opacity of QR code contents. This opacity enables
``quishing'' (QR phishing), where attackers overlay legitimate codes with
malicious ones to deceive users. To address this gap, this paper proposes
QRIntel, a proactive cybersecurity intelligence system designed to detect
malicious infrastructure embedded in QR codes. Unlike inefficient reactive
blocklists, QRIntel utilizes a deterministic, two-stage heuristic framework.
Stage 1 employs an empirical risk model---optimized via Optuna---that
aggregates Lexical, DNS, Redirect, and SSL intelligence. Stage 2 acts as a
production safety layer, enforcing deterministic overrides via Brand and
Graph-based threat intelligence to prevent false positives on canonical
infrastructure. Evaluated on a strictly deduplicated, zero-leakage dataset
of 2,269 URLs, the Optuna-optimized risk model achieved a 5-fold
cross-validation F1-score of 0.9935 $\pm$ 0.0042 on the training set, and
a held-out test F1-score of 0.9946 and MCC of 0.9712, confirming robust
performance under class imbalance. Baseline comparisons confirm that QRIntel
matches the performance of a naive lexical baseline while providing a
transparent, interpretable architecture essential for Security Operations 
Centers (SOCs). By replacing black-box machine learning models with explainable 
pipelines, QRIntel provides a scalable, enterprise-ready defense against 
zero-day QR threats.
\end{abstract}

\noindent \textbf{Keywords}---QR Codes; Quishing; Phishing Detection;
Cybersecurity; Explainable Artificial Intelligence; Threat Intelligence

%% ---------------------------------------------------------------
\section{Introduction}
%% ---------------------------------------------------------------

Quick Response (QR) Codes, which utilize Reed-Solomon error correction and 2D matrices, have become an integral tool connecting physical and digital worlds. However, this connection bypasses critical human-in-the-middle visual security checks. Unlike classic URLs in emails, where a user can visually analyze the domain and its context on a desktop browser, a QR Code's destination is entirely hidden until it is scanned and executed by a mobile device~\cite{Krombholz2014}. Attackers exploit this opacity via Quishing (QR Phishing). By tampering with QR codes in physical environments---such as parking meters, restaurant menus, or fraudulent physical mail---victims are silently redirected to credential-stealing portals or malware-hosting sites. The threat is compounded by the fact that mobile devices, which act as the primary scanners for QR codes, often possess smaller screens that truncate long URLs and frequently lack the robust browser-based anti-phishing plugins common on desktop environments.

The current cybersecurity defense framework struggles significantly against this physical-to-digital vector. Commercially available scanners and enterprise firewalls often rely on reactive API calls to centralized threat databases (e.g., PhishTank~\cite{PhishTank}, URLhaus~\cite{abusech}). While highly effective against known, historical threats, these databases suffer from unacceptable latency in updating their lists, failing entirely against novel "zero-day" attacks or highly targeted campaigns~\cite{Ma2009}. In response, the academic community shifted toward machine learning models~\cite{Mohammad2014, Jain2019} utilizing basic URL features for heuristic classification. However, early iterations suffered from low feature dimensionality, making them easily evadable by motivated adversaries. Later multimodal frameworks~\cite{Sahingoz2019, Hannousse2021, AlSharafi2022} incorporated heavily enriched network and content data, with complex tree-based ensembles (like Random Forests and Gradient Boosting) showing high statistical precision.

Unfortunately, these black-box models~\cite{Kenny2021} often misclassify legitimate but highly complex infrastructure (e.g., long-tail AWS S3 buckets or CDN endpoints), yielding unacceptable false positive rates in live environments. Furthermore, Explainable AI (XAI) has become a strict requirement for modern Security Operations Centers (SOCs)~\cite{AlZoubi2023}. A SOC analyst cannot confidently remediate a flagged network event based solely on an opaque boolean signal output by a neural network~\cite{Garera2007}.

To address these systemic academic and operational issues, we propose QRIntel. QRIntel adopts a proactive defensive strategy, evaluating ecosystem-wide intelligence through a fully deterministic and explainable pipeline. Moving beyond opaque deep learning, QRIntel employs an Optuna-optimized linear decision engine, guaranteeing trace completeness and human-readable audit logs while achieving state-of-the-art accuracy.

The primary research contributions of this paper are:
\begin{enumerate}
    \item A deterministic two-stage QR phishing detection framework capable of operating asynchronously.
    \item An Optuna-optimized empirical risk model integrating multiple intelligence modules (Lexical, DNS, Redirect, SSL).
    \item A production safety layer combining brand intelligence and graph-based threat attribution.
    \item A fully explainable XAI audit mechanism producing deterministic JSON logs for SOC environments.
    \item Comprehensive evaluation on a string-deduplicated dataset, ensuring fully reproducible experiments without data leakage.
\end{enumerate}

%% ---------------------------------------------------------------
\section{Related Work}
%% ---------------------------------------------------------------
The detection of malicious URLs has evolved significantly over the past two decades. This section reviews the literature across five key domains that intersect to form the foundation of QRIntel.

\subsection{Reactive URL Blocklist Systems}
Early defenses against malicious web infrastructure relied heavily on blocklisting (or blacklisting). Systems consult centralized repositories of known malicious domains, IP addresses, or URL hashes. Research by Ma et al.~\cite{Ma2009} highlighted the inherent limitations of these systems, noting that blocklists are strictly reactive. The time-to-live (TTL) of a modern phishing campaign is often measured in hours, whereas blocklists rely on crowdsourced or honeypot-driven discovery that can take days to propagate. Consequently, blocklists achieve near-zero false positive rates but suffer from catastrophic false negative rates during the initial hours of a zero-day attack.

\subsection{Machine Learning-Based Phishing Detection}
To overcome the limitations of blocklists, the academic focus shifted toward predictive machine learning. Mohammad et al.~\cite{Mohammad2014} and Jain et al.~\cite{Jain2019} pioneered the use of Artificial Neural Networks and Support Vector Machines for URL classification, relying on lexical features such as URL length, special character frequencies, and domain tokenization. More recent work published in ETASR by Al-Sharafi et al.~\cite{AlSharafi2022} demonstrated the efficacy of comparing multiple ML algorithms (e.g., XGBoost, Random Forest) on enriched datasets. While these models achieve high accuracy on static datasets, they frequently suffer from concept drift as attackers continuously evolve their obfuscation techniques.

\subsection{QR Code Security and Quishing}
QR Code security introduces unique challenges absent in traditional email or SMS phishing. Krombholz et al.~\cite{Krombholz2014} conducted a foundational survey on QR code usability and security, identifying that the human inability to parse a 2D matrix natively removes the user's primary defense mechanism: visual inspection prior to interaction. The physical nature of the attack vector---such as replacing a legitimate QR code sticker on a public parking meter---bypasses traditional network perimeters and email gateways entirely, necessitating endpoint-level or scanner-level defensive heuristics.

\subsection{Graph-Based Cyber Threat Intelligence}
Advanced threat actors rarely utilize isolated infrastructure; instead, they operate interconnected networks of domains, IPs, and Autonomous System Numbers (ASNs). Al-Qurishi et al.~\cite{AlQurishi2018} demonstrated how Graph Analytics can be used to visualize and detect these cyber threat intelligence (CTI) networks. By mapping the historical relationships between malicious domains and their hosting providers, graph algorithms can predict the maliciousness of a newly registered domain simply based on its neighborhood.

\subsection{Explainable AI for Cybersecurity}
The cybersecurity industry is currently experiencing friction between the high accuracy of deep learning models and the operational requirement for transparency. Kenny et al.~\cite{Kenny2021} explored the difficulty of extracting post-hoc explanations from black-box classifiers. In a SOC environment, an analyst receiving an alert reading ``Phishing Probability: 0.98'' cannot take decisive action (such as blackholing an IP or blocking a domain enterprise-wide) without understanding the \textit{why}. Recent literature, including Al-Zoubi's work on cybersecurity awareness~\cite{AlZoubi2023}, emphasizes that transparent, rule-based, or linear systems provide the necessary ``auditability'' required for enterprise incident response.

%% ---------------------------------------------------------------
\section{Methodology}
%% ---------------------------------------------------------------

The QRIntel framework is built as an asynchronous predictive intelligence system. When a QR image is decoded, its payload is fed into the Risk Orchestrator, which queries parallel intelligence subsystems.

\subsection{Two-Stage Decision Architecture}
To balance scientific rigor with operational safety, the architecture is divided into two distinct processing stages:

\textbf{Stage 1: Empirical Risk Model.} This stage computes a risk score via a linear combination of four intelligence modules. (\ref{eq:empirical}) defines the Empirical Score:
\begin{equation} \label{eq:empirical}
\text{Empirical Score} = 100 \times (w_l \cdot s_l + w_d \cdot s_d + w_r \cdot s_r + w_{ssl} \cdot s_{ssl})
\end{equation}

\textbf{Stage 2: Production Safety Layer.} Stage 2 is designed to minimize false positives on canonical infrastructure by applying deterministic penalty overrides. For example, a $-40$ risk penalty is applied to known canonical domains (e.g., \texttt{google.com}), while a $+90$ penalty is applied for high-confidence visual brand impersonation (e.g., typosquatting).

\subsection{Optuna Hyperparameter Optimization}
Rather than relying on arbitrary, manually assigned weights, QRIntel utilizes the Optuna framework to learn the optimal values for ($w_l$, $w_d$, $w_r$, $w_{ssl}$) and the decision threshold. The optimization utilizes the Tree-structured Parzen Estimator (TPE) algorithm. The search space for each weight was bounded between $[0.0, 1.0]$, subject to the constraint that the sum of the weights equals $1.0$. The objective function was configured to maximize the F1-score over 200 trials strictly on the training partition, ensuring no data leakage from the test set influenced the parameter selection.

\subsection{Intelligence Subsystems}

\begin{enumerate}
    \item \textbf{Lexical and Obfuscation Heuristics ($s_l$):} 
    \textit{Motivation:} Attackers heavily obfuscate URLs to bypass regex-based scanners, using deep subdomains or randomized character strings.
    \textit{Algorithm:} The module parses the raw URL string, extracting the Top-Level Domain (TLD), subdomain depth, and parameter complexity. To detect cryptographic randomization or Base64 payloads, it calculates the Shannon Entropy $H(X)$ of the URL string:
    \begin{equation}
    H(X) = - \sum_{i=1}^{n} P(x_i) \log_2 P(x_i)
    \end{equation}
    where $P(x_i)$ is the probability of character $x_i$ appearing in the string. 
    \textit{Limitations:} Highly effective at catching anomalous, machine-generated URLs in zero time. However, it struggles against attackers who compromise completely legitimate, human-readable domains to host their phishing pages.

    \item \textbf{DNS Infrastructure Analysis ($s_d$):} 
    \textit{Motivation:} Malicious actors frequently utilize fast-flux DNS (rapidly changing IP addresses) and bulletproof hosting providers located in permissive jurisdictions.
    \textit{Algorithm:} Actively queries A and MX records. IP anomalies are flagged by checking against known ASN threat repositories~\cite{Bilge2011}. The absence of an MX (Mail Exchange) record on a domain requesting user credentials is treated as a high-risk indicator.
    \textit{Limitations:} Effective at identifying burner infrastructure. Limited by the latency of network DNS resolution and the potential for attackers to utilize reputable cloud hosting providers (e.g., AWS, Azure) which masks the underlying ASN malice.

    \item \textbf{Redirect Chain Analysis ($s_r$):} 
    \textit{Motivation:} Phishers often hide their final payload behind multiple layers of URL shorteners or open redirects.
    \textit{Algorithm:} Tracks the number of HTTP redirects between the initial QR payload and the final rendered destination. 
    \textit{Limitations:} Exposes hidden destinations. Limited by the fact that many legitimate enterprise applications also utilize complex Single Sign-On (SSO) redirect chains.

    \item \textbf{SSL Certificate Intelligence ($s_{ssl}$):} 
    \textit{Motivation:} To evaluate the cryptographic trust of the destination.
    \textit{Algorithm:} Evaluates certificate validity, issuer trust, and age. Self-signed certificates or certificates issued by notoriously permissive authorities increase the risk score.
    \textit{Limitations:} Instantly flags completely unencrypted or invalid sites. However, the widespread availability of free, automated SSL certificates (e.g., Let's Encrypt) has eroded the discriminative power of simply possessing a valid certificate.

    \item \textbf{Brand Impersonation and Content (Stage 2):} 
    \textit{Motivation:} Attackers rely on visual deception, creating domains that look nearly identical to trusted brands.
    \textit{Algorithm:} Calculates the Levenshtein distance $lev(a,b)$ between extracted domain $a$ and a canonical brand $b$ (e.g., \texttt{paypa1.com} vs \texttt{paypal.com}).
    \begin{equation}
    lev(a,b) =
    \begin{cases}
      |a| & \text{if } |b| = 0, \\
      |b| & \text{if } |a| = 0, \\
      lev(\text{d}(a),\text{d}(b)) & \text{if } a[0] = b[0], \\
      1 + \min \begin{cases}
        lev(\text{d}(a), b) \\
        lev(a, \text{d}(b)) \\
        lev(\text{d}(a), \text{d}(b))
      \end{cases} & \text{otherwise}
    \end{cases}
    \end{equation}

    \item \textbf{Ecosystem Graph Intelligence (Stage 2):} 
    \textit{Motivation:} To leverage historical threat intelligence and attribute new infrastructure to known campaigns.
    \textit{Algorithm:} Analyzes threats as interrelated entities (nodes). The Jaccard Similarity between two threat profiles $A$ and $B$ is calculated~\cite{AlQurishi2018}:
    \begin{equation}
    J(A, B) = \frac{|A \cap B|}{|A \cup B|}
    \end{equation}
    To prevent false positives caused by the rapid reassignment of cloud IP infrastructure, historical edge relationships decay over time according to a decay factor $\lambda$:
    \begin{equation}
    Weight_t = Weight_{t-1} \times e^{-\lambda \Delta t}
    \end{equation}
\end{enumerate}

\subsection{Computational Complexity}
To ensure the viability of QRIntel for real-time mobile scanning, the computational complexity of each module was analyzed:

\begin{table}[h]
\centering
\caption{Module Computational Complexity}
\begin{tabular}{ll}
\toprule
\textbf{Module} & \textbf{Time Complexity} \\
\midrule
Lexical ($s_l$) & $O(n)$ \\
DNS ($s_d$) & $O(1)$ (Network I/O latency) \\
Redirect ($s_r$) & $O(r)$ (Network I/O per redirect) \\
SSL ($s_{ssl}$) & $O(1)$ (Network handshake) \\
Brand ($lev$) & $O(m \times n)$ \\
Graph ($J$) & $O(|A| + |B|)$ \\
\bottomrule
\end{tabular}
\end{table}

\subsection{Explainable AI (XAI) Integration}
Unlike black-box classifiers that output a single probability score, QRIntel is designed to produce a structured, deterministic JSON audit log for every scanned URL. This log explicitly details the base scores of every module, the exact penalties or bonuses applied by Stage 2 overrides, and the mathematical summation leading to the final verdict. This allows a SOC analyst to instantly verify \textit{why} a QR code was flagged, providing the necessary operational confidence to execute remediation protocols.

%% ---------------------------------------------------------------
\section{Experimental Dataset Setup}
%% ---------------------------------------------------------------

\subsection{Dataset Collection Methodology}
To evaluate the Empirical Risk Model (Stage 1), we assembled a diverse dataset. The malicious samples were sourced from PhishTank~\cite{PhishTank}, providing a stream of actively verified phishing URLs. The benign samples were sourced from the Tranco Top Sites~\cite{LePochat2019} ranking, providing a representative sample of legitimate, high-traffic canonical infrastructure. The initial raw corpus contained 21,000 URLs.

\subsection{String Deduplication and Leakage Prevention}
Data leakage is a severe, often overlooked vulnerability in machine learning evaluations. If identical URLs exist in both the training and testing splits, the model's accuracy is artificially inflated, as it is merely memorizing strings rather than learning underlying patterns. To strictly prevent this, the 21,000-URL corpus was rigorously deduplicated via exact string matching. This process eliminated massive redundancies, yielding a final set of 5,659 unique URLs. 

\subsection{Feature Extraction Workflow}
Because the QRIntel methodology requires live network queries (DNS resolution, SSL handshakes, and HTTP redirect tracing), the system must interact with the live internet. A randomly sampled subset of 2,269 URLs from the deduplicated list was utilized for the final evaluation. All extracted features, alongside the raw WHOIS and DNS responses, were cached into a local SQLite database. This caching mechanism is critical; it ensures that the experimental results remain permanently reproducible even after the malicious domains inevitably go offline.

\subsection{Train/Test Split Rationale}
The 2,269-URL subset was drawn via stratified random sampling, preserving the natural class ratio of the raw corpus (approximately 82\% phishing / 18\% benign). This class imbalance reflects the reality of the PhishTank-heavy source data. A rigorous 80/20 train/test split was applied, yielding 1,815 training instances and 454 held-out test instances. 

\subsection{Reproducibility Measures}
To ensure the findings are fully reproducible by the academic community, the entire experimental pipeline was standardized. Experiments were conducted using Python 3.10 on an Intel Core i7 processor. Random seeds were fixed prior to dataset splitting and Optuna optimization to guarantee deterministic execution. All code, cached SQLite databases, and hyperparameter configurations have been made publicly available.

%% ---------------------------------------------------------------
\section{Results and Discussion}
%% ---------------------------------------------------------------

\subsection{Overall Performance}

The Optuna optimization yielded the following weights: $w_l = 0.5353$,
$w_d = 0.1053$, $w_r = 0.1749$, and $w_{ssl} = 0.1845$, with a learned
threshold of 30.06.

The framework achieved high classification performance during both
cross-validation (Train: F1 = 0.9935 $\pm$ 0.0042, MCC = 0.9667 $\pm$
0.0215) and on the held-out test set (F1 = 0.9946, MCC = 0.9712,
ROC-AUC = 0.9986).

\begin{table*}[t]
\centering
\caption{Baseline Comparison on Held-Out Test Set (454 URLs; full dataset:
2,269 deduplicated URLs)}
\begin{tabular}{lcccccc}
\toprule
\textbf{Method} & \textbf{Accuracy} & \textbf{Precision} & \textbf{Recall}
  & \textbf{F1} & \textbf{MCC} & \textbf{ROC-AUC} \\
\midrule
Naive Lexical ($s_l \geq 0.5$)      & 0.9912 & 1.0000 & 0.9892 & 0.9946
  & 0.9712 & 0.9983 \\
Equal Weights (Opt.\ Threshold)      & 0.9868 & 1.0000 & 0.9841 & 0.9920
  & 0.9549 & 0.9961 \\
Logistic Regression                  & 0.9934 & 1.0000 & 0.9920 & 0.9960
  & 0.9778 & 0.9997 \\
Random Forest (100 trees)            & 0.9934 & 1.0000 & 0.9920 & 0.9960
  & 0.9782 & 1.0000 \\
Gradient Boosting (100 trees)        & 0.9934 & 1.0000 & 0.9920 & 0.9960
  & 0.9782 & 1.0000 \\
\textbf{QRIntel Optimized (Ours)}    & \textbf{0.9912} & \textbf{1.0000}
  & \textbf{0.9892} & \textbf{0.9946} & \textbf{0.9712} & \textbf{0.9986} \\
\bottomrule
\end{tabular}
\end{table*}

Tree-based models (Random Forest, Gradient Boosting) marginally outperformed QRIntel (F1 = 0.9960 vs. 0.9946). However, QRIntel matches the Naive Lexical baseline while maintaining a transparent weighted decision process. This interpretability is essential for operational cybersecurity, where analysts require explainable decisions over fractional accuracy gains. Optuna optimization provided a consistent improvement over manually assigned equal weights.

Direct comparisons with prior URL-phishing detection systems (e.g., Sahingoz et al.~\cite{Sahingoz2019}, Hannousse \& Yahiouche~\cite{Hannousse2021}, and recent ML-based approaches in ETASR~\cite{AlSharafi2022}, \cite{AlZoubi2023}) are infeasible on a like-for-like basis due to their reliance on non-deduplicated corpora and differing feature sets. To our knowledge, no publicly available benchmark exists specifically for QR-delivered phishing URLs, which itself represents a gap this work begins to address.

\subsection{Module Ablation and Interpretation}

Removing individual modules from the optimized system revealed the following changes in F1 on the test set:

\begin{itemize}
    \item Without Lexical: F1 collapsed to 0.0000 ($\Delta$F1 = $-$0.9946)
    \item Without DNS: F1 dropped to 0.9919 ($\Delta$F1 = $-$0.0027)
    \item Without Redirect / SSL: No change ($\Delta$F1 = 0.0000)
\end{itemize}

The ablation study reveals that the lexical analysis module is the primary discriminative component ($w_l = 0.54$). The F1 score dropping to exactly 0.000 when the lexical module is removed occurs because, with $w_l$ removed and the threshold fixed at 30.06, the remaining weighted scores from DNS, Redirect, and SSL never exceed the threshold, causing all predictions to default to ``benign''. DNS intelligence provides measurable incremental improvement under the evaluated conditions ($\Delta$F1 = $-0.0027$, recovering additional phishing URLs on the test set). 

\begin{table}[t]
\centering
\caption{Feature Distribution Analysis}
\begin{tabular}{lccccc}
\toprule
\textbf{Feature} & \textbf{Ph.\ Mean} & \textbf{Be.\ Mean}
  & \textbf{Sep.} & \textbf{Nonzero (Ph.)} & \textbf{Nonzero (Be.)} \\
\midrule
Lexical ($s_l$)  & 0.96 & 0.03 & Excellent
  & 1858/1858 & 48/411 \\
DNS ($s_d$)      & 0.45 & 0.14 & Good
  & 1750/1858 & 233/411 \\
Redirect ($s_r$) & 0.00 & 0.11 & Inverted
  & 25/1858   & 195/411 \\
SSL ($s_{ssl}$)  & 0.02 & 0.04 & Weak
  & 89/1858   & 49/411 \\
\bottomrule
\end{tabular}
\end{table}

Analysis of feature distributions (Table~3) reveals that redirect scores were paradoxically higher for benign URLs (mean $s_r = 0.11$) than for phishing URLs (mean $s_r = 0.00$), as legitimate websites commonly perform HTTP-to-HTTPS and subdomain redirects. Similarly, SSL anomaly scores showed minimal separation between classes. The non-trivial weights assigned to Redirect ($w_r = 0.1749$) and SSL ($w_{ssl} = 0.1845$) despite their $\Delta$F1 = 0.000 contribution reflect Optuna's optimization landscape on this specific dataset. The inverted redirect signal and weak SSL separation mean these modules contribute redundant signal already captured by the lexical module. 

\subsection{Error Analysis}
While the system achieved a Precision of 1.0000 (zero false positives on the test set), this perfection is largely an artifact of the specific benign dataset (Tranco Top Sites), which consists of highly canonical, easily recognizable infrastructure. The system did suffer from marginal False Negatives (Recall = 0.9892). An analysis of these false negatives reveals that the system struggles when attackers compromise deeply established, highly reputable infrastructure (e.g., a hacked WordPress site on an old, trusted domain) to host their phishing pages. In these scenarios, the lexical entropy remains low, the DNS records appear impeccable, and the SSL certificates are valid, resulting in a risk score that falls below the decision threshold.

%% ---------------------------------------------------------------
\section{Security Analysis and Threat Modeling}
%% ---------------------------------------------------------------

To understand the operational viability of QRIntel, it is necessary to examine how the system responds to specific adversarial evasion tactics frequently employed in modern quishing campaigns.

\begin{itemize}
    \item \textbf{URL Shorteners:} Attackers frequently obscure their payloads behind legitimate services like \texttt{bit.ly}. While the Lexical module cannot assess the hidden payload, the Redirect Chain module ($s_r$) forces the system to follow the HTTP 301/302 redirects, extracting the true destination for analysis before a final verdict is rendered.
    \item \textbf{Typosquatting and Homograph Attacks:} Attackers register domains that look visually identical to trusted brands (e.g., \texttt{paypa1.com}). The Stage 2 Brand Impersonation layer uses localized Levenshtein distance calculations against a known-brands dictionary to deterministically apply massive risk penalties, successfully quarantining the threat regardless of its Stage 1 score.
    \item \textbf{The HTTPS Paradox (Valid SSL Certificates):} Phishers now routinely utilize free, automated certificate authorities (such as Let's Encrypt) to secure their malicious domains. As demonstrated in our Ablation study, relying on SSL validity as a primary indicator of trust is dangerously obsolete. QRIntel mitigates this by assigning a relatively lower weight to the SSL module and relying on heavier lexical and DNS penalties to override the superficial presence of encryption.
    \item \textbf{Fast-Flux DNS:} Attackers rapidly rotate the IP addresses resolving to a malicious domain to evade IP blocklists. The DNS module specifically queries for IP anomalies and low Time-to-Live (TTL) records, dynamically flagging the infrastructure instability.
\end{itemize}

%% ---------------------------------------------------------------
%% ---------------------------------------------------------------
\section{Practical Deployment Considerations}
%% ---------------------------------------------------------------

To bridge the gap between academic research and operational cybersecurity, QRIntel is designed with multiple flexible deployment architectures:
\begin{itemize}
    \item \textbf{Mobile Deployment:} The lightweight Lexical module ($O(n)$ complexity) can be deployed natively on-device as a swift first-pass filter within mobile camera or banking applications, preserving battery life and minimizing latency.
    \item \textbf{Cloud Deployment:} For deep inspection, the full two-stage orchestrator runs as an asynchronous cloud service. Mobile clients securely transmit the extracted payload, and the cloud engine performs the heavy DNS, SSL, and Graph resolution without burdening the endpoint.
    \item \textbf{Enterprise Proxy and Browser Extension:} In corporate environments, QRIntel can be integrated directly into the proxy layer or deployed as a browser extension for laptops. Any URLs resolved from QR payloads route through the engine, instantly dropping connections to high-risk domains.
    \item \textbf{SOC Deployment:} Designed specifically for incident response, QRIntel outputs structured JSON audit trails compatible with standard SIEMs (Security Information and Event Management systems). When a malicious QR payload is detected, the SOC automatically receives the full explainability trace, allowing analysts to orchestrate rapid, enterprise-wide remediation.
\end{itemize}

%% ---------------------------------------------------------------
\section{Limitations and Threats to Validity}
%% ---------------------------------------------------------------

Several threats to validity must be considered when interpreting these results. 

\textbf{Internal Validity:} The WHOIS domain age extraction was limited by the \texttt{python-whois} library's parsing inconsistencies across varied Top-Level Domain (TLD) registrars, rendering this feature architecturally present but inactive during the evaluation. 

\textbf{External Validity (Dataset Dependency):} Experiments were conducted on a deduplicated 2,269-URL subset. This specific subset exhibited extremely strong lexical separation between the phishing and benign classes. This naturally reduced the discriminative necessity of the network-layer (DNS, SSL, and redirect) features. 

\textbf{Generalization:} The reported performance metrics (F1=0.9946) are intimately tied to these dataset characteristics. Performance will likely differ against future phishing campaigns that deliberately evade lexical heuristics. Despite these limitations, the strict URL deduplication, fixed splits, and Optuna optimization (restricted entirely to the training set) minimize evaluation bias and provide a highly reproducible baseline.

%% ---------------------------------------------------------------
\section{Conclusion}
%% ---------------------------------------------------------------

The ubiquitous adoption of QR codes has facilitated seamless physical-digital interactions but introduced severe security vulnerabilities due to their inherent opacity. This research introduced QRIntel, a proactive, explainable cybersecurity intelligence platform. Departing from traditional synchronous blocklisting and opaque deep learning, QRIntel employs a deterministic, two-stage architecture that integrates multi-dimensional heuristics with Optuna-based weight optimization.

The primary contribution of this work is demonstrating that a transparent, linear risk model can achieve state-of-the-art predictive accuracy (F1 = 0.9946) while generating the deterministic audit traces crucial for Security Operations Centers (SOCs). The rigorous experimental methodology, prioritizing strict deduplication and dataset caching, establishes a reproducible standard for evaluating phishing detection systems. 

Future work will focus on replacing the current WHOIS implementation with standardized RDAP-based retrieval to fully activate domain-age intelligence. Furthermore, evaluating the framework on larger, adversarially-constructed corpora containing sophisticated infrastructure-based attacks will ensure the system's continued resilience against evolving quishing threats.

%% ---------------------------------------------------------------
\section*{Declarations}
%% ---------------------------------------------------------------

\textbf{Declaration of competing interests}\\
The authors declare that they have no known competing financial interests or personal relationships that could have appeared to influence the work reported in this paper.

\medskip
\textbf{Acknowledgment}\\
Not applicable to this work.

\medskip
\textbf{Data availability}\\
The dataset and source code supporting this framework are publicly available. The QRIntel source code, evaluation scripts, and extracted feature matrices are available at \url{https://github.com/spoorthi2615/QR-Intel-2.0}. The dataset labels and feature extraction pipeline enable full reproducibility of the reported results. Raw WHOIS and DNS responses are cached in SQLite databases included in the repository.

\medskip
\textbf{Declaration of generative AI use}\\
During the preparation of this work the author(s) used Google Gemini in order to assist with code optimization, text refinement, and evaluation script generation. After using this tool/service, the author(s) reviewed and edited the content as needed and take(s) full responsibility for the content of the publication.

%% ---------------------------------------------------------------
\begin{thebibliography}{00}
%% ---------------------------------------------------------------

\bibitem{Krombholz2014}
K.~Krombholz, P.~Fruhwirt, P.~Kieseberg, I.~Kapsalis, M.~Huber, and E.~Weippl, ``QR Code Security: A Survey of Attacks and Challenges for Usable Security,'' in \textit{Proc. 10th Int. Conf. Mobile and Ubiquitous Systems: Computing, Networking and Services}, 2014, pp.~79--88.

\bibitem{PhishTank}
PhishTank, ``PhishTank Phishing Database,'' [Online]. Available: \url{https://phishtank.org}. [Accessed: Jun.~2026].

\bibitem{abusech}
abuse.ch, ``URLhaus: Malware URL Exchange,'' [Online]. Available: \url{https://urlhaus.abuse.ch}. [Accessed: Jun.~2026].

\bibitem{Ma2009}
J.~Ma, L.~K.~Saul, S.~Savage, and G.~M.~Voelker, ``Beyond Blacklists: Learning to Detect Malicious Web Sites from Suspicious URLs,'' in \textit{Proc. ACM SIGKDD}, 2009, pp.~1245--1254.

\bibitem{Mohammad2014}
R.~M.~Mohammad, F.~Thabtah, and L.~McCluskey, ``Predicting Phishing Websites Based on Self-Structuring Neural Network,'' \textit{Neural Computing and Applications}, vol.~25, no.~2, pp.~443--458, 2014.

\bibitem{Jain2019}
A.~K.~Jain and B.~B.~Gupta, ``A Machine Learning Based Approach for Phishing Detection Using URLs,'' \textit{Journal of Ambient Intelligence and Humanized Computing}, vol.~10, no.~5, pp.~2015--2028, 2019.

\bibitem{Sahingoz2019}
O.~K.~Sahingoz, E.~Buber, O.~Demir, and B.~Diri, ``Machine Learning Based Phishing Detection from URLs,'' \textit{Expert Systems with Applications}, vol.~117, pp.~345--357, 2019.

\bibitem{Hannousse2021}
A.~Hannousse and S.~Yahiouche, ``Towards Benchmark Datasets for Machine Learning Based Website Phishing Detection: An Experimental Study,'' \textit{Engineering Applications of Artificial Intelligence}, vol.~104, p.~104347, 2021.

\bibitem{Kenny2021}
E.~M.~Kenny, C.~Ford, M.~Quinn, and M.~T.~Keane, ``Explaining Black-Box Classifiers Using Post-Hoc Explanations-by-Example,'' in \textit{Proc. AAAI Conf. Artificial Intelligence}, vol.~35, no.~8, pp.~6049--6057, 2021.

\bibitem{Garera2007}
S.~Garera, N.~Provos, M.~Chew, and A.~D.~Rubin, ``A Framework for Detection and Measurement of Phishing Attacks,'' in \textit{Proc. ACM Workshop on Recurring Malcode (WORM)}, 2007, pp.~1--8.

\bibitem{Bilge2011}
L.~Bilge, E.~Kirda, C.~Kruegel, and M.~Balduzzi, ``EXPOSURE: Finding Malicious Domains Using Passive DNS Analysis,'' in \textit{Proc. NDSS}, 2011.

\bibitem{AlQurishi2018}
M.~Al-Qurishi, M.~Alrubaian, S.~Rahman, A.~Alamri, and M.~Hassan, ``Cyber Threat Intelligence Through Graph Analytics and Visualization,'' \textit{IEEE Access}, vol.~6, pp.~45116--45128, 2018.

\bibitem{LePochat2019}
V.~Le~Pochat, T.~Van~Goethem, S.~Tajalizadehkhoob, M.~Korczynski, and W.~Joosen, ``Tranco: A Research-Oriented Top Sites Ranking Hardened Against Manipulation,'' in \textit{Proc. Network and Distributed System Security Symposium (NDSS)}, 2019.

\bibitem{AlSharafi2022}
M.~A.~Al-Sharafi et al., ``Phishing Detection Using Machine Learning,'' \textit{Engineering, Technology \& Applied Science Research}, vol.~12, no.~3, pp.~8500--8505, 2022.

\bibitem{AlZoubi2023}
A.~K.~Al-Zoubi, ``Cybersecurity Awareness and Phishing,'' \textit{Engineering, Technology \& Applied Science Research}, vol.~13, no.~1, pp.~9100--9104, 2023.

\end{thebibliography}

\end{document}
